% Reconstructed after the environment replacement of 20 November 2026.
% This is a fresh complete transcription, not a byte-for-byte recovery.
\subsection{Uniform quadratic representation growth on local punctures}

Throughout this subsection, \(R\) is a complete excellent normal local
domain of dimension two, of equal characteristic \(p>0\), with an
algebraically closed coefficient field \(k\) and maximal ideal
\(\mathfrak m\). We put \(U=\operatorname{Spec}(R)\setminus\{\mathfrak m\}\).
A relative maximal Cohen--Macaulay module over
\(R_A=R\otimes_k A\), for a finite-dimensional local Artinian \(k\)-algebra \(A\)
with residue field \(k\), means
a finite \(A\)-flat module whose closed fibre is maximal
Cohen--Macaulay. All tensor products with \(A\) here are already complete,
since \(A\) has finite length over \(k\).

\begin{lemma}[Proved: relative Noether-different parameters and prescribed jets]
\label{D:relative-Noether-different-jets}
There is a system of parameters \(x,y\) of \(R\) with the following
properties. Put \(I_b=(x^b,y^b)R_A\).
For every relative maximal Cohen--Macaulay module \(M_A\),
multiplication by \(x\) and by \(y\) factors through finite free
\(R_A\)-modules. Consequently, if \(N_A\) is any finite \(R_A\)-module
on which \(x,y\) is a regular sequence, every homomorphism
\[
 M_A/I_bM_A\longrightarrow N_A/I_bN_A,\qquad b\geq2,
\]
lifts to a homomorphism \(M_A\to N_A\) agreeing with the prescribed
map modulo \(I_{b-1}\).
The target is not required to be flat over \(A\).

For \(H=\operatorname{Hom}_{R_A}(M_A,N_A)\) and
\(K_c=\operatorname{Hom}_{R_A}(M_A,I_cN_A)\), we also have
\[
 I_cH\subset K_c
 \subset (x^c,y^c,x^{c-1}y^{c-1})H
 \subset I_{c-1}H \qquad(c\geq2).
 \tag{L1}
\]
\end{lemma}

\begin{proof}
A normal local surface is Cohen--Macaulay and is regular away from its
closed point. The sum of its Noether differents is therefore
\(\mathfrak m\)-primary
\cite[Theorem 1.1(1), definition on p.~298 and proof on p.~300]
{DaoTakahashi2015B}. Choose parameters \(x,y\) in that sum.
Only finitely many Noether normalizations \(S_i\cong k[[u_i,v_i]]\)
are needed to express them as sums of elements of the corresponding
differents. The ring \(R\) is finite free over each \(S_i\).

We spell out the factorization because its relative form is essential.
For \(S_i\subset R\), let \(\mu:R\otimes_{S_i}R\to R\) be multiplication,
and choose
\(\eta=\sum_\nu u_\nu\otimes v_\nu\) in
\(\operatorname{Ann}_{R\otimes_{S_i}R}(\ker\mu)\).
If \(z=\mu(\eta)\), the maps
\[
 M_A\longrightarrow R_A\otimes_{S_{i,A}}M_A
       \longrightarrow M_A,\qquad
 m\longmapsto\sum_\nu u_\nu\otimes v_\nu m,\quad
 r\otimes m\longmapsto rm
\]
are \(R_A\)-linear and their composite is \(z\,\operatorname{Id}\).
Here \(S_{i,A}=S_i\otimes_k A\).
The middle module is finite free over \(R_A\): indeed the closed fibre
of \(M_A\) is free over \(S_i\), and a lifted basis gives
\(S_{i,A}^{a}\to M_A\). Its cokernel vanishes by nilpotent Nakayama.
Flatness over \(A\) makes its kernel reduce to zero modulo the maximal
ideal of \(A\); nilpotent Nakayama kills that kernel as well.
Taking direct sums of the displayed factorizations proves the assertion
for \(x,y\). It also shows that these elements annihilate
\(\operatorname{Ext}^1_{R_A}(M_A,N)\) for every target \(N\).

Let \(\bar f\) be the prescribed \(b\)-jet. Reduce first modulo \(x^b\).
The factorization of \(y\,\operatorname{Id}_{M_A}\) reduces to a
factorization through a free \(R_A/(x^b)\)-module. The obstruction to
lifting \(y\bar f\) across multiplication by \(y^b\) therefore vanishes.
We obtain
\[
 g:M_A/x^bM_A\longrightarrow N_A/x^bN_A,\qquad
 g\bmod y^b=y\bar f.
\]
The image of \(g\) is contained in \(y(N_A/x^bN_A)\).
Division by the nonzerodivisor \(y\) gives a map \(h\) agreeing with
\(\bar f\) modulo \(y^{b-1}\). Now the \(x\)-annihilation of
\(\operatorname{Ext}^1_{R_A}(M_A,N_A)\) lifts \(xh\) to a map
\(M_A\to N_A\). Its image is in \(xN_A\); divide by \(x\).
The resulting map agrees with \(\bar f\) modulo
\((x^{b-1},y^{b-1})\), as asserted.

For (L1), the Koszul presentation
\[
 0\longrightarrow N_A\longrightarrow N_A^{\,2}
       \longrightarrow I_cN_A\longrightarrow0
\]
and the same annihilation give
\(xK_c,yK_c\subset I_cH\).
The sequence \(x,y\) is regular on \(H\). For example, if \(yh=xh'\),
regularity of \(y\) on \(N_A/xN_A\) implies that every value of \(h\)
lies in \(xN_A\), and division defines a homomorphism \(h/x\).
The usual two-parameter colon calculation, valid on every module
with this regular sequence, gives
\[
 (I_cH:(x,y))=(x^c,y^c,x^{c-1}y^{c-1})H.
\]
This proves all the inclusions.
\end{proof}

\begin{lemma}[Proved: exact realization of relative Artin module data]
\label{D:relative-Artin-module-realization}
Set \(S=k[[x,y]]\), let \(d=\operatorname{rank}_S R\), and fix \(L\geq3\).
Put \(J_L=(x^L,y^L)\subset S\).
An action of
\[
 R_A/I_L^2 \quad\hbox{on}\quad
 (S_A/J_L^2)^{dr}
\]
extending its given \(S_A/J_L^2\)-module structure has a reduction \(T_A\)
modulo \(I_L\) which is the exact reduction of a relative maximal
Cohen--Macaulay \(R_A\)-module \(M_A\).
Its closed fibre has rank \(r\) over \(R\), and it is free of rank \(dr\)
over \(S_A\). Conversely every rank-\(r\) maximal Cohen--Macaulay
\(R\)-module supplies such data.

Over \(k\), the isomorphism class of \(M\) is determined by
\(M/I_LM\). At fixed \(L\), the action data above are described by
\(O(r^2)\) scalar variables and \(O(r^2)\) polynomial equations of
bounded degree, with constants depending on \(R,L\), not on \(r\).
\end{lemma}

\begin{proof}
Write \(C_2=S_A/J_L^2\), \(C_1=S_A/J_L\), and let \(W_A\) be the given
\(C_2\)-free module. A free \(C_2\)-resolution of \(C_1\), tensored
with the \(C_2\)-free algebra \(R_A/I_L^2\), is a free resolution
of \(R_A/I_L\). Thus
\[
 \operatorname{Tor}^{R_A/I_L^2}_i(R_A/I_L,W_A)
   =\operatorname{Tor}^{C_2}_i(C_1,W_A)=0\quad(i>0).
\]
This is a lifting in the full Tor sense, not just a lifting of an
underlying presentation.
Apply \cite[Proposition 3.6, pp.~18--19, including its proof]
{AuslanderDingSolberg1993B} to the regular sequence \(x^L,y^L\).
It says that \(T_A\) is a direct summand of \(N_A/I_LN_A\), where
\[
 N_A=\Omega^2_{R_A}(T_A)\oplus P_A
\]
and \(P_A\) is finite free. Since \(T_A\) is \(A\)-free, the two
syzygies are \(A\)-flat. The closed fibre of \(N_A\) is maximal
Cohen--Macaulay. The basis-lifting argument in the preceding lemma
therefore makes \(N_A\) free over every \(S_{i,A}\).

Let \(\bar e\) be the projector onto \(T_A\) in \(N_A/I_LN_A\).
The prescribed-map lemma lifts it to \(f\in\operatorname{End}_{R_A}(N_A)\)
agreeing modulo \(I_{L-1}\).
Hence \(f^2-f\) lies in the closed two-sided ideal
\[
 K=\{h:h(N_A)\subset I_{L-1}N_A\}.
\]
By (L1), \(K\subset I_{L-2}\operatorname{End}_{R_A}(N_A)\).
It is topologically nilpotent in this complete endomorphism algebra.
The elementary idempotent Newton iteration therefore produces an
idempotent \(e\) with \(e-f\in K\).
Modulo \(I_L\), the projectors \(e'\) and \(\bar e\) differ by the
nilpotent ideal of maps whose images lie in
\(I_{L-1}N_A/I_LN_A\). Idempotents congruent modulo a nilpotent ideal
are conjugate: explicitly
\[
 u=e'\bar e+(1-e')(1-\bar e)
\]
is congruent to \(1\) modulo that ideal and satisfies
\(e'u=u\bar e\). Consequently
\(M_A=eN_A\) has reduction isomorphic to \(T_A\) exactly, not merely
to a smaller truncation. It is a direct summand of a module free over
every \(S_{i,A}\), hence is free over those local rings. Its
\(S_A\)-rank is \(dr\), by reduction, and its closed-fibre \(R\)-rank
is \(r\).

For uniqueness over \(k\), lift a prescribed isomorphism
\(M/I_LM\cong M'/I_LM'\) by the preceding lemma.
The lift is an isomorphism modulo \(\mathfrak m\), hence is surjective.
The two modules are \(S\)-free of the same rank, so it is an isomorphism.
The converse is reduction of the \(S\)-free module \(M\).

Finally choose an \(S\)-basis \(z_1,\ldots,z_d\) of \(R\).
The action is a fixed finite list of \(dr\) by \(dr\) matrices over
the fixed finite-dimensional algebra \(S/J_L^2\).
The multiplication table of \(R\), the unit equation, and the
\(S/J_L^2\)-linearity give \(O(r^2)\) equations of degree at most two
in \(O(r^2)\) scalar coordinates. No determinant inversion of
rank-dependent degree is needed.
\end{proof}


\begin{lemma}[Proved: quadratic witnesses for cyclic Frobenius data]
\label{D:quadratic-cyclic-Frobenius-witnesses}
There are constants \(c,D\), depending only on \(R\), such that for
every \(r,f\geq1\) there is a closed affine \(k\)-scheme \(V_{r,f}\)
in at most \(cfr^2\) scalar coordinates, defined by at most \(cfr^2\)
equations of degrees at most \(D\), with the following properties.
Every \(k\)-point gives a cyclic sequence of rank-\(r\) reflexive modules
and isomorphisms on \(U\)
\[
 F_U^*(M_j|_U)\ \cong\ M_{j+1}|_U
 \qquad(j\in\mathbf Z/f\mathbf Z).
 \tag{L2}
\]
Every such sequence occurs. More strongly, over every local Artinian
\(k\)-algebra \(A\), every point of \(V_{r,f}(A)\) gives actual
relative maximal Cohen--Macaulay modules and such a unit Frobenius
cycle on \(U_A\). The actions at a fixed Artin depth are exactly the
reductions of these modules.
\end{lemma}

\begin{proof}
Fix once and for all
\[
 B=3,\qquad L=pB=3p,\qquad n=p(B-1)=2p,\qquad
 \mathfrak a=(x^p,y^p).
 \tag{L3}
\]
Then \(I_n\subset\mathfrak m\mathfrak a\).
For every reflexive \(M\), the natural map
\(F_R^*M\to F_R^{[*]}M\) has cokernel annihilated by \(\mathfrak a\).
Here is a proof of the uniformity. Apply Frobenius to a factorization
of \(x\operatorname{Id}_M\) through a finite free module.
It factors \(x^p\operatorname{Id}_{F^*M}\) through a free module.
The map from \(F^*M\) to that free module extends to the reflexive hull,
because on \(U\) it is a map between vector bundles and maps into
a free module extend across the closed point. The resulting identity,
first on \(U\) and then by reflexivity, shows
\(x^pF^{[*]}M\subset\operatorname{im}(F^*M)\).
The same argument applies to \(y\).

At each stage take the action matrices of
Lemma~\ref{D:relative-Artin-module-realization} at depth \(I_L^2\).
Their reductions \(T_j\) modulo \(I_L\) are free of rank \(dr\) over
\(S/J_L\).
An arrow
\[
 a_j:T_j/I_BT_j\longrightarrow T_{j+1}
\]
is required to be \(p\)-semilinear. Its matrix \(A_j\) has entries
in \(S/J_L\), and the equations for the fixed \(S\)-basis \(z_i\)
of \(R\) are
\[
 A_j\,\rho_j(z_i)^{[p]}=\rho_{j+1}(z_i)^p A_j.
 \tag{L4}
\]
The notation on the left means entrywise Frobenius; the power on
the right is matrix multiplication and represents multiplication
by \(z_i^p\). The \(S\)-semilinearity is built into the matrix
description. In particular the source is killed by \(I_B\) because
\(I_B^{[p]}=I_L\) kills the target.

For each \(j\) and \(b\in\{x^p,y^p\}\), introduce matrices
\(X_{i,j,b}\) over \(S/J_n\), and impose
\[
 \rho_{j+1}(b)=
   \sum_{i=1}^d \rho_{j+1}(z_i)\,A_jX_{i,j,b}
       \pmod {I_n}.
 \tag{L5}
\]
This says that the image of multiplication by \(\mathfrak a\)
is contained, at depth \(n\), in the \(R\)-span of the semilinear
arrow. The \(z_i\) span \(R\) over \(S\), so these matrices are
precisely inclusion witnesses, with no omitted coefficients.
Equations (L4) have degree at most \(p+1\), and (L5) degree at
most three. All coefficient rings and the number \(d\) are fixed.
Together with the quadratic action equations, there are
\(O(fr^2)\) variables and equations of degree bounded independently
of \(r,f\).

We must prove that every witness, including a nonreduced one, is
effective. Let \(A\) be local Artinian.
Realize its module actions by the preceding lemma, obtaining
\(M_{j,A}\). View the arrow as an \(R_A\)-linear map into
\(F_*M_{j+1,A}/I_BF_*M_{j+1,A}\).
The latter quotient has underlying module \(M_{j+1,A}/I_LM_{j+1,A}\).
The target \(F_*M_{j+1,A}\) is finite and \(x,y\)-regular:
\(R_A\) is \(F\)-finite, and \(x,y\) act on it as \(x^p,y^p\).
It need not be \(A\)-flat. The prescribed-map lemma nevertheless
lifts the arrow to an actual semilinear map, agreeing in the
underlying target modulo \(I_n\).

Let \(C\) be the cokernel of its linearization
\(F^*M_{j,A}\to M_{j+1,A}\).
The lifted equations (L5) imply
\[
 \mathfrak a C\subset I_nC\subset\mathfrak m\mathfrak a C.
\]
Nakayama gives \(\mathfrak a C=0\). On \(U_A\),
\(\mathfrak a\) is the unit ideal, so the map is surjective there.
Both modules are vector bundles of rank \(r\) on \(U_A\), hence
the map is an isomorphism. This proves (L2) over \(A\).
In particular no flatness of the Frobenius-restriction target
was used.

Conversely, a unit cycle gives reflexive modules at all stages.
The uniform defect calculation above supplies (L5), and reduction
of its actual maps supplies (L4). Since each \(M_j\) is \(S\)-free,
it supplies the required action at depth \(I_L^2\).
Thus every desired cycle has a witness.
\end{proof}

\begin{lemma}[Proved: uniqueness of a unit structure on a fixed module]
\label{D:local-complete-Lang}
Let \(q=p^f\), and let \(M\) be a reflexive \(R\)-module with a unit
\(q\)-Frobenius structure on \(U\).
Any two such structures on the same underlying module are isomorphic.
\end{lemma}

\begin{proof}
Put \(\mathcal E=\operatorname{End}_R(M)\), a finite complete
\(R\)-algebra. A chosen unit structure induces a multiplicative
additive \(q\)-semilinear map
\(\tau:\mathcal E\to\mathcal E\), by Frobenius pullback and conjugation.
Maps on \(U\) extend to \(M\), so this is a map on the whole algebra.
It satisfies
\[
 \tau(\mathfrak m^a\mathcal E)\subset
       \mathfrak m^{qa}\mathcal E.
 \tag{L6}
\]
The finite-dimensional \(k\)-algebra
\(\overline{\mathcal E}=\mathcal E/\mathfrak m\mathcal E\)
has a connected smooth unit group: it is the determinant-nonzero
open subset of its underlying affine vector space, using its
left regular representation. The induced \(\bar\tau\) has zero
differential. Every twisted-conjugacy orbit
\[
 g\longmapsto g\,a\,\bar\tau(g)^{-1}
\]
is therefore open, since its orbit map has invertible differential.
Connectedness forces a single orbit. Hence the equation
\(a=g\bar\tau(g)^{-1}\) is solvable modulo \(\mathfrak m\).

The ideal \(\mathfrak m\mathcal E\) lies in the Jacobson radical,
so a solution modulo it lifts to a unit \(g_0\in\mathcal E^\times\).
With the orientation \(g=a\tau(g)\), iterate
\(g_{j+1}=a\tau(g_j)\). By (L6), successive errors tend to zero
\(\mathfrak m\)-adically with exponents multiplied by \(q\).
The complete algebra supplies a unit limit solving the equation.
This conjugates the two unit structures. We have used Lang's
argument only on the finite-dimensional residue algebra, not on
the infinite-dimensional \(k\)-space \(\mathcal E\).
\end{proof}

\begin{theorem}[Proved: uniform quadratic local representation growth]
\label{D:uniform-local-representation-growth}
There exists \(C_R>0\) such that, for every \(r,f\geq1\) and \(q=p^f\),
\[
 \#\bigl(
 \operatorname{Hom}_{\mathrm{cts}}(\pi_1(U),\operatorname{GL}_r(\mathbf F_q))
                  /\operatorname{GL}_r(\mathbf F_q)\bigr)
       \leq q^{C_Rr^2}.
 \tag{L7}
\]
The constant is independent of both the rank and the finite coefficient
field.
\end{theorem}

\begin{proof}
Finite \(\mathbf F_q\)-local systems on \(U\) are equivalent to
vector bundles equipped with a unit \(q\)-Frobenius structure.
For completeness, after trivializing the bundle the fixed-vector
equation is \(v=A v^{[q]}\), with \(A\) invertible. Multiplication
by \(A^{-1}\) makes the leading monomials \(v_i^q\), so the
fixed-vector scheme is finite locally free of degree \(q^r\).
Its Jacobian is invertible, hence it is finite etale.
Over an algebraically closed geometric fibre its fixed vectors
form an \(r\)-dimensional \(\mathbf F_q\)-space spanning the
ambient vector space: apply Frobenius to a shortest normalized
linear relation and subtract to see that the coefficients belong
to \(\mathbf F_q\). The invertible fixed frames therefore form
the required finite etale \(\operatorname{GL}_r(\mathbf F_q)\)-torsor.
Conversely the associated bundle of such a torsor has its evident
unit structure. Its vector bundle extends uniquely to a finite
reflexive module on the normal local surface.

A \(q\)-unit structure can be written as an \(f\)-cycle of
\(p\)-Frobenius structures by taking the successive reflexive
Frobenius pullbacks. Thus every representation class is represented
on \(V=V_{r,f}\) of the preceding lemma.
We prove that each representation class is an open and closed
part of the underlying topological space of \(V\).

The key assertion is stronger than constancy on geometric points.
Fix a closed point \(v\in V(k)\), and consider any point over a
local Artinian \(k\)-algebra \(A\) reducing to \(v\).
Its genuine relative unit cycle consists of \(M_{j,A}\).
Choose a multiple \(e\) of \(f\) so large that the \(p^e\)-power
map kills the maximal ideal of \(A\).
Absolute Frobenius on \(R_A\) then factors through the closed
fibre followed by its \(p^e\)-Frobenius.
Iterating the unit cycle identifies
\[
 M_{j,A}|_{U_A}\cong M_{j,0}|_U\otimes_k A.
 \tag{L8}
\]
Here we use the iterated special-fibre unit structure to identify
the pulled-back special bundle with \(M_{j,0}|_U\).
Both sides satisfy Hartogs extension across \(V(\mathfrak m R_A)\):
the regular sequence \(x,y\) gives vanishing of local cohomology
in degrees zero and one. Hence (L8) extends to an isomorphism
of the complete relative modules.

Only their canonical reductions modulo \(I_L\) are asserted to
be constant. The chosen action modulo \(I_L^2\), the auxiliary
inclusion witnesses, and the chosen arrow lifts can still vary.
This distinction is crucial.

Let \(\mathcal T\) be the affine scheme of tuples of actions of
\(R/I_L\) on the fixed free \(S/J_L\)-module of rank \(dr\).
Reduction of our witness data gives a morphism
\(\pi:V\to\mathcal T\).
The algebraic group
\[
 \mathcal G=
 \operatorname{Res}_{(S/J_L)/k}
          \operatorname{GL}_{dr}^{\,f}
\]
acts by changing the bases at the \(f\) stages.
The preceding Artinian assertion says that every deformation of
\(v\), at every finite Artin order, has its image in the ordinary
\(\mathcal G\)-orbit of \(\pi(v)\).
These ordinary orbits are locally closed.
One can also see their scheme structure directly: the stabilizer
is the unit group of a finite-dimensional endomorphism algebra,
hence is smooth, and the orbit is the corresponding homogeneous
locally closed subscheme.

Choose an ambient open neighbourhood of \(\pi(v)\) in which its
orbit is closed. The ideal of that orbit, pulled back to the
Noetherian local ring \(\mathcal O_{V,v}\), vanishes in every
quotient by a power of its maximal ideal, by (L8).
Krull intersection makes it zero in \(\mathcal O_{V,v}\).
Its finite set of generators therefore vanishes on a neighbourhood
of \(v\). This neighbourhood maps into the orbit.

At closed points, belonging to this orbit means that the canonical
module tuples modulo \(I_L\) are isomorphic. The exact realization
lemma lifts the isomorphisms to the reflexive modules.
Lemma~\ref{D:local-complete-Lang} says that the underlying periodic
module determines its unit \(q\)-structure up to isomorphism.
Thus the neighbourhood just constructed consists of one
representation class.

The class neighbourhoods for distinct classes are disjoint:
a nonempty intersection in a finite-type \(k\)-scheme would have
a closed point. Their union contains all closed points, and hence
all of \(V\), since \(V\) is Jacobson. Quasicompactness gives a
finite clopen partition by classes. In particular the number of
classes is at most the number of irreducible components of \(V\).

Finally, a closed affine scheme cut out by \(m\) equations of
degree at most \(D\) has at most \(D^m\) irreducible components.
Indeed, start with affine space and intersect with the equations
one at a time. A component contained in a hypersurface is retained;
otherwise the hypersurface degree formula bounds the sum of
the degrees of the new components by \(D\) times its old degree.
Induction proves the bound, also when dimensions vary.
Here \(m\leq cfr^2\), so the number is at most \(D^{cfr^2}\).
Increasing \(C_R\) gives (L7).
\end{proof}

\begin{remark}[Scope of the counting theorem]
\label{D:uniform-local-growth-scope}
The theorem bounds isomorphism classes, equivalently conjugacy classes
of representations. It does not assert a bound on the number of raw
homomorphisms without the additional conjugation factor, nor a
uniform bound on generators of the full fundamental group.
The proof uses exact liftability of all Artinian witnesses and
openness of their class fibres. A finite-dimensional parameter
space, or mere determinacy of the special-fibre jet, would not
by itself justify the component count.
\end{remark}

\begin{corollary}[Proved: the bound for higher-dimensional local punctures]
\label{D:uniform-local-growth-all-dimensions}
Let \(A\) be a complete excellent normal local domain of equal
characteristic \(p>0\), of dimension at least two, with algebraically
closed residue field. Its punctured spectrum has a constant \(C_A\)
for which the bound \(q^{C_A r^2}\) holds for every \(q=p^f\) and \(r\).
\end{corollary}

\begin{proof}
The local surface reduction of
\ref{D:local-surface-reduction} gives a surjection
\[
 \left(\mathop{*}_{i=1}^{t}\pi_1(U_i)\right)
                *\widehat F_e\longrightarrow
       \pi_1(\operatorname{Spec}(A)\setminus\{\mathfrak m_A\}),
\]
where the \(U_i\) are finitely many complete normal local surface
punctures with the same algebraically closed residue field.
For each factor choose a representative of each representation
class. Simultaneous conjugacy leaves at most
\(\lvert\operatorname{GL}_r(\mathbf F_q)\rvert^{e+t-1}\)
choices for relative conjugations and for the free generators
(with the harmless empty-factor case treated separately).
Multiplying the finitely many surface bounds and using
\(\lvert\operatorname{GL}_r(\mathbf F_q)\rvert\leq q^{r^2}\)
proves the assertion. Pullback along the displayed surjection
is injective on representation classes.
\end{proof}


